\section{评估：效率指标的误导性}
\label{sec:eval}

\S\ref{sec:reality} 已给出核心的实测结果。本节从评估方法论的角度做一层反思：为什么常规的效率指标不仅无法发现前述退化，反而会掩盖它？并给出应当采用的有效性指标。

\subsection{效率指标全部“通过”}

若以自主代理系统常用的效率指标衡量这次部署，结论会是全面正面的（表~\ref{tab:misleading}）：规则命中率 $\rho=100\%$，意味着几乎零 LLM 成本；决策延迟极低（规则匹配为毫秒级）；无崩溃、无未捕获异常，系统连续运行三天。任何只看这一组指标的评估都会判定“系统高效、稳定、经济”。

\begin{table}[t]
  \centering\small
  \caption{效率指标（真实值）与其误导性解读。}
  \label{tab:misleading}
  \begin{tabular}{@{}lll@{}}
    \toprule
    指标 & 实测值 & 表面解读 vs.\ 真相 \\
    \midrule
    规则命中率 $\rho$ & 100\% & “省成本” / 实为智能缺位 \\
    LLM 调用次数 & 0 & “高效” / 推理层被饿死 \\
    决策延迟 & 毫秒级 & “响应快” / 只是正则匹配 \\
    崩溃次数 & 0 & “稳定” / 稳定地空转 \\
    \bottomrule
  \end{tabular}
\end{table}

这正是 \S\ref{sec:rootcause} 所述“指标伪装”的具体体现：效率指标度量的是“系统做一件事有多快多省”，而完全不关心“这件事是否有意义”。

\subsection{应采用的有效性指标}

要暴露退化，必须引入度量\emph{任务是否真正推进}的指标。我们定义并在本次数据上测量如下三项：

\textbf{有效动作率。} 定义为“带来可观测任务进展的动作数 / 总动作数”。以“注入了真实语义内容”为宽松代理指标，本次部署中携带真实任务指令的动作仅 1 次，有效动作率约 $1/46\approx 2\%$。即便把 5 次 CLI 重启也算作“有意图的动作”，上限也不过 $6/46\approx13\%$。

\textbf{自主交付成功率。} 定义为“真正产出内容的自主子任务数 / 尝试数”。据 \S\ref{sec:reality} 发现三，RalphEngine 的 headless 执行\textbf{无一产出内容}，该指标为 $0\%$。

\textbf{空转比例。} 定义为“未带来进展的重复动作占比”。40 次“继续”中绝大多数属于此类，空转比例接近 $87\%$。

三项有效性指标（约 2\% 有效动作、0\% 自主交付、87\% 空转）与三项效率指标（$\rho=100\%$、0 次 LLM、0 崩溃）构成尖锐对比。同一次部署，在两套指标下得出完全相反的结论。这一对比本身就是本文的一个方法论贡献：\textbf{评估自主代理系统时，效率指标必须与有效性指标并列，否则“高效的空转”会被误报为成功。}

\subsection{有效性威胁}

\textbf{样本规模：} 本文基于单次、单机、46 次决策的部署，规模有限，不足以支撑统计性结论；其价值在于提供一个有代码根因支撑的、可复现的负面案例，而非普适的量化率。\textbf{代理指标近似：} “有效动作率”以“是否注入真实语义内容”近似“是否推进任务”，可能低估或高估真实有效性；更严格的度量需要对任务状态做独立标注。\textbf{外部有效性：} 退化的具体触发依赖于所用 CLI 的界面与规则实现，其他配置下的表现可能不同，但 \S\ref{sec:rootcause} 所抽象的“廉价层垄断”反模式不依赖具体实现。
